\section{Theoretical Background}
\label{sec:theory}

\subsection{Bang-Bang Control and Limit Cycles}
\label{sec:theory:bangbang}

A Bang-Bang (two-point) controller with hysteresis~$d$ generates the
manipulated variable signal
\begin{equation}
    u(t) =
    \begin{cases}
        U_{\max} & \text{if } e(t) > +d, \\
        U_{\min} & \text{if } e(t) < -d, \\
        u(t^-)   & \text{otherwise},
    \end{cases}
    \label{eq:bangbang}
\end{equation}
where $e(t) = r - y(t)$ is the control error, $r$ is the setpoint, $y(t)$ is
the plant output, and $u(t^-)$ denotes the previous value (memory).

\Cref{fig:block_closed_loop} illustrates the closed-loop structure.
For a stable plant $G(s)$ with sufficient gain and phase delay, this loop
converges to a periodic \emph{limit-cycle oscillation}~\cite{Atherton1975}.
Let $\Tu$ be the oscillation period and $\Au$ the peak-to-peak amplitude of
$y(t)$ in steady limit-cycle operation.

\begin{figure}[htbp]
    \centering
    \begin{tikzpicture}[
        >=Stealth, auto,
        block/.style={draw, rectangle, minimum height=0.75cm,
                      minimum width=1.5cm, align=center, font=\small},
        sum/.style={draw, circle, inner sep=2pt, font=\small},
        node distance=0.5cm and 1.0cm,
    ]
        \node[sum]                          (sum)  {$\Sigma$};
        \node[block, right=1.2cm of sum]    (bb)   {Bang-Bang\\[-2pt]$d,\,U_{\min},U_{\max}$};
        \node[block, right=1.2cm of bb]     (plant){$G(s)$};
        \coordinate[right=0.9cm of plant]   (yjunc);
        \coordinate[below=0.7cm of yjunc]   (fb);
        \coordinate[below=0.7cm of sum]     (fbsum);

        \draw[->] ++(-1.2,0) node[left]{$r$} -- node[above,font=\small]{} (sum);
        \draw[->] (sum) -- node[above,font=\small]{$e$} (bb);
        \draw[->] (bb)  -- node[above,font=\small]{$u$} (plant);
        \draw[->] (plant) -- node[above,font=\small]{$y$} (yjunc)
                           -- ++(0.5,0) node[right,font=\small]{$y$};
        \draw[-]  (yjunc) -- (fb) -- (fbsum);
        \draw[->] (fbsum) -- (sum.south);
        \node[below left=1pt and 1pt of sum, font=\small]{$-$};
    \end{tikzpicture}
    \caption{Closed-loop block diagram: Bang-Bang controller drives the plant
             $G(s)$ into a sustained limit-cycle oscillation around the
             setpoint~$r$.}
    \label{fig:block_closed_loop}
\end{figure}

\subsection{Describing Function Analysis}
\label{sec:theory:df}

The describing function (DF) of a relay with hysteresis~$d$ and amplitude~$M$
(where $M = (U_{\max} - U_{\min})/2$) is~\cite{Atherton1982}

\begin{equation}
    N(A) = \frac{4M}{\pi A}
           \left( \sqrt{1 - \left(\frac{d}{A}\right)^2}
                  - j \frac{d}{A} \right),
    \quad A > d,
    \label{eq:df}
\end{equation}
where $A = \Au/2$ is the single-sided oscillation amplitude of the output.

The limit-cycle frequency~$\omega_u = 2\pi/\Tu$ satisfies the harmonic balance
condition
\begin{equation}
    G(j\omega_u) = -\frac{1}{N(A)}.
    \label{eq:hb}
\end{equation}

For the special case of zero hysteresis ($d = 0$), \eqref{eq:df} reduces to
the ideal relay describing function $N(A) = 4M/(\pi A)$, and
\eqref{eq:hb} becomes equivalent to the Ziegler--Nichols ultimate gain
condition $|G(j\omega_u)| = \pi A / (4M)$.

\subsection{FOPDT Plant Model}
\label{sec:theory:fopdt}

The First-Order-Plus-Dead-Time (FOPDT) transfer function
\begin{equation}
    G(s) = \frac{K e^{-Ls}}{Ts + 1}
    \label{eq:fopdt}
\end{equation}
is a standard approximation for a large class of industrial processes.
It is characterised by three parameters: process gain~$K$, time
constant~$T$, and dead time~$L$.
\Cref{fig:fopdt_block} shows the internal structure of this model.

\begin{figure}[htbp]
    \centering
    \begin{tikzpicture}[
        >=Stealth, auto,
        block/.style={draw, rectangle, minimum height=0.75cm,
                      minimum width=1.7cm, align=center, font=\small},
        node distance=0.4cm and 1.0cm,
    ]
        \node[block] (delay) {$e^{-Ls}$};
        \node[block, right=1.2cm of delay] (fo) {$\dfrac{K}{Ts+1}$};

        \draw[->] ++(-1.5,0) node[left,font=\small]{$u(t)$} -- (delay);
        \draw[->] (delay) -- (fo);
        \draw[->] (fo) -- ++(1.5,0) node[right,font=\small]{$y(t)$};
    \end{tikzpicture}
    \caption{Internal block diagram of the FOPDT model~\eqref{eq:fopdt}:
             the dead-time element delays the input, which then drives a
             first-order lag with gain~$K$ and time constant~$T$.}
    \label{fig:fopdt_block}
\end{figure}

\subsubsection{Parameter Identification from Limit Cycle}

Evaluating \eqref{eq:fopdt} at $s = j\omega_u$ gives
\begin{align}
    |G(j\omega_u)| &= \frac{K}{\sqrt{1 + (\omega_u T)^2}},
    \label{eq:gain_cond} \\
    \angle G(j\omega_u) &= -\arctan(\omega_u T) - \omega_u L = -\pi.
    \label{eq:phase_cond}
\end{align}

The describing function provides two equations (gain and phase of
$-1/N(A)$):
\begin{align}
    |G(j\omega_u)| &= \frac{\pi A}{4M}
                      \frac{1}{\sqrt{1 - (d/A)^2}},
    \label{eq:gain_from_df} \\
    \angle G(j\omega_u) &= -\pi + \arcsin\!\left(\frac{d}{A}\right).
    \label{eq:phase_from_df}
\end{align}

Combining \eqref{eq:phase_cond} and~\eqref{eq:phase_from_df} yields
\begin{equation}
    \arctan(\omega_u T) + \omega_u L
    = \pi - \arcsin\!\left(\frac{d}{A}\right).
    \label{eq:phase_id}
\end{equation}

With $\omega_u$ known from the measured period $\Tu$, \eqref{eq:gain_cond}, \eqref{eq:gain_from_df} and~\eqref{eq:phase_id}
form a system that can be solved for $K$, $T$, and $L$ if one additional
assumption is introduced.
A common engineering choice is to adopt the ratio $L/T$ from a
characteristic of the oscillation waveform (e.g., the symmetry of the
on/off intervals) or to set $L = \Tu/(2\pi) \cdot \phi_L$ based on the
phase margin~\cite{Seborg2011}.

\begin{remark}
    For processes where the FOPDT approximation is poor (e.g., resonant
    systems), higher-order describing function methods or relay feedback
    with filtered output should be considered.
\end{remark}

\begin{remark}
    While \eqref{eq:gain_cond}, \eqref{eq:gain_from_df} and~\eqref{eq:phase_id} characterise the
    limit cycle and guarantee the existence and uniqueness of the FOPDT
    parameters $K$, $T$, $L$, the KZV does not solve this nonlinear system
    directly. Instead, Phase~3 (\cref{sec:method:phase3}) employs a
    transient-slope method to identify $\hat{L}$, $\hat{T}$, $\hat{K}$
    sequentially from the time-domain signal. This approach is more robust
    in practice because it avoids the numerical sensitivity of the nonlinear
    DF system while remaining theoretically justified by the DF equations.
\end{remark}

\subsection{IMC-Based PID Tuning}
\label{sec:theory:imc}

Given the FOPDT model~\eqref{eq:fopdt}, the IMC-PID tuning
rules~\cite{Rivera1986} yield

\begin{align}
    \Kp &= \frac{1}{K} \cdot \frac{T}{\lambda + L}, \label{eq:Kp} \\
    \Ti  &= T,                                         \label{eq:Ti} \\
    \Td  &= \frac{L}{2},                               \label{eq:Td}
\end{align}
where $\lambda > 0$ is the desired closed-loop time constant and serves as a
robustness tuning knob: larger $\lambda$ gives a more sluggish but more
robust response.

A recommended default is $\lambda = L$ (i.e., equal to the estimated dead
time). This choice, consistent with the guideline $\lambda \geq 0.5L$
recommended in~\cite{Rivera1986}, provides a closed-loop bandwidth
approximately equal to the reciprocal of twice the dead time, yielding a
good trade-off between speed and robustness for typical industrial processes.

\begin{remark}
    The modified ZN formulas for FOPDT (see, e.g.,~\cite{Aastrom2006}) are
    an alternative and are compared against the IMC rules in
    \cref{sec:results}.
\end{remark}
